\subsection{Verificación de Factibilidad por Invariante de Clases}

\graybold{Preguntas guía:}

\begin{itemize}
\item ¿Los elementos están agrupados en clases (colores, tipos) y solo se permiten intercambios DENTRO de la misma clase?
\item ¿Se pregunta si una configuración objetivo es alcanzable, en vez de contarlas o construirlas?
\item ¿La clase de cada posición es invariante (no cambia con los intercambios permitidos), y solo importa qué elemento termina en cada posición?
\end{itemize}

\graybold{Utilidad:} Verificar si una configuración objetivo es alcanzable cuando las únicas operaciones permitidas son intercambios dentro de una misma clase (color, grupo). Como los intercambios dentro de una clase permiten CUALQUIER permutación entre sus posiciones, el problema se reduce a comparar, para cada posición, si la clase requerida coincide con la clase disponible.

\graybold{Complejidad:} O(n log n) por el ordenamiento.

\graybold{Fundamento teórico:} Intercambiar repetidamente elementos de una misma clase genera el grupo simétrico completo sobre las posiciones de esa clase: cualquier permutación interna es alcanzable. Por lo tanto, la clase (color) de cada posición es un invariante que ningún intercambio permitido puede romper. La configuración ordenada es alcanzable si y solo si, al ordenar los elementos por su valor objetivo, la clase que le corresponde a cada uno coincide con la clase original de esa posición.

\graybold{Ejercicio aplicado}

Una secuencia de n bloques, cada uno con un número distinto (1..n) y un color, se puede reordenar intercambiando bloques del mismo color. Determinar si es posible dejar la secuencia ordenada ascendentemente por número.

\graybold{I/O:} n ≤ 10\textasciicircum{}5 → lectura línea a línea (patrón 1) es aceptable dado que el costo dominante es el ordenamiento, no la lectura.

\graybold{Código solución}

\begin{lstlisting}[language=Python]
import sys
input = sys.stdin.buffer.readline
 
def solve():
    n, k = map(int, input().split())
    color_at_position = [0] * n     # color de cada posicion (invariante)
    blocks = []                     # (numero, color) de cada bloque
    for i in range(n):
        num, c = map(int, input().split())
        color_at_position[i] = c
        blocks.append((num, c))
 
    blocks.sort()                   # ordena por numero: donde deberia quedar cada uno
    # si el color del bloque que llega a la posicion i coincide con el color
    # que YA tenia esa posicion, es alcanzable (el color nunca cambia de lugar)
    ok = all(blocks[i][1] == color_at_position[i] for i in range(n))
    print('Y' if ok else 'N')
 
solve()
\end{lstlisting}
